\fancyhead[L]{\small BEITRÄGE ZUR THEORIE}
\SourceLeaf{080}
\begin{center}
{\large\bfseries BEITRÄGE ZUR THEORIE\par}
\medskip
{\Large\bfseries DER ALGEBRAISCHEN GLEICHUNGEN.\par}
\medskip
\rule{0.18\linewidth}{0.4pt}
\end{center}

Es werden in dieser Denkschrift zwei verschiedene die algebraischen
Gleichungen betreffende Gegenstände behandelt. Zuerst stelle ich den
vor funfzig Jahren von mir gegebenen Beweis des Grundlehrsatzes der
Theorie der algebraischen Gleichungen in einer veränderten Gestalt
und mit erheblichen Zusätzen auf. Der zweite Theil ist einer speciellen
Behandlung der algebraischen Gleichungen mit drei Gliedern gewidmet,
und enthält Methoden, nicht bloss die reellen, sondern auch die
imaginären Wurzeln solcher Gleichungen mit Leichtigkeit zu bestimmen.

\begin{center}
\rule{0.18\linewidth}{0.4pt}\par
\medskip
\textbf{ERSTE ABTHEILUNG.}
\end{center}

Die im Jahre 1799 erschienene Denkschrift, \emph{Demonstratio nova
theorematis, omnem functionem algebraicam rationalem integram unius
variabilis in factores reales primi vel secundi gradus resolvi posse},
hatte einen doppelten Zweck, nämlich erstens, zu zeigen, dass
sämmtliche bis dahin versuchte Beweise dieses wichtigsten Lehrsatzes
der Theorie der algebraischen Gleichungen ungenügend und illusorisch
sind, und zweitens, einen neuen vollkommen strengen Beweis zu geben.
Es ist unnöthig, auf den erstern Gegenstand noch einmal zurückzukommen.
Dem dort gegebenen neuen Beweise habe ich selbst später noch zwei
andere folgen lassen, und ein vierter ist zuerst von
\textsc{Cauchy} aufgestellt. Diese vier Beweise \SourceLeaf{081}beruhen alle auf eben so vielen verschiedenen Grundlagen,
aber darin kommen sie alle überein, dass durch jeden derselben
zunächst nur das Vorhandensein \emph{Eines} Factors der betreffenden
Function erwiesen wird. Der Strenge der Beweise thut dies allerdings
keinen Eintrag: denn es ist klar, dass wenn von der vorgegebenen
Function dieser eine Factor abgelöset wird, eine ähnliche Function
von niederer Ordnung zurückbleibt, auf welche der Lehrsatz aufs
neue angewandt werden kann, und dass durch Wiederholung des
Verfahrens zuletzt eine vollständige Zerlegung der ursprünglichen
Function in Factoren der bezeichneten Art hervorgehen wird.
Indessen gewinnt ohne Zweifel jede Beweisführung eine höhere
Vollendung, wenn nachgewiesen wird, dass sie geeignet ist, das
Vorhandensein der sämmtlichen Factoren unmittelbar anschaulich
zu machen. Dass der erste Beweis in diesem Fall ist, habe ich
bereits in der gedachten Denkschrift angedeutet (Art. 23), ohne
es dort weiter auszuführen: dies soll jetzt ergänzt werden, und
ich benutze zugleich diese Gelegenheit, die Hauptmomente des ganzen
Beweises in einer abgeänderten und, wie ich glaube, eine
vergrösserte Klarheit darbietenden Gestalt zu wiederholen.
Was dabei die äussere Einkleidung des Lehrsatzes selbst betrifft,
so war die 1799 gebrauchte, dass die Function
$x^n+Ax^{n-1}+Bx^{n-2}+\mathrm{u.\,s.\,w.}$ sich in reelle
Factoren erster oder zweiter Ordnung zerlegen lässt, damals
deshalb gewählt, weil alle Einmischung imaginärer Grössen vermieden
werden sollte. Gegenwärtig, wo der Begriff der complexen Grössen
jedermann geläufig ist, scheint es angemessener, jene Form fahren
zu lassen und den Satz so auszusprechen, dass jene Function sich
in $n$ \emph{einfache} Factoren zerlegen lasse, wo dann die
constanten Theile dieser Factoren nicht eben reelle Grössen zu sein
brauchen, sondern für dieselben auch jede complexen Werthe
zulässig sein müssen. Bei dieser Einkleidung gewinnt selbst der
Satz noch an Allgemeinheit, weil dann die Beschränkung auf reelle
Werthe auch bei den Coëfficienten $A,B$ u.\,s.\,w. nicht
vorausgesetzt zu werden braucht, vielmehr jedwede Werthe für
dieselben zulässig bleiben.

\FifthArticle{1}
Wir betrachten demnach die Function der unbestimmten Grösse $x$
\[
x^n+Ax^{n-1}+Bx^{n-2}+\mathrm{u.\,s.\,w.}+Mx+N=X
\]
wo $A,B\,.\,.\,.\,.\,M,N$ bestimmte reelle oder imaginäre Coëfficienten
vorstellen.
